
\section{Free Church Meetings}

Editorial article in “Folkebladet” for March 16, 1898.

Both Dr. (J. A. S.) and Pastor Kildahl seem to think that it is far better to have the association government over the congregations than that the congregations, entirely free and independent, should decide all their own matters for themselves. And when these men hold this opinion, they surely have the right to maintain it.

But this does not give them the right to present matters as though the Free Church meetings had the same authority as the old association meetings, though they were constituted in an entirely different manner.

Let us come to a clear understanding of this matter; for otherwise we do one another wrong without intending it.

The Free Church meetings are constituted differently from the associations’ annual meetings, and therefore they have absolutely no power or authority over a single congregation.

With respect to the congregation’s right, the Free Church has laid down this principle in its work, that no congregation shall or may be bound by any decision other than its own.

Now one may certainly hold the opinion, if one wishes, that it is foolish to work according to such a principle, and that the congregations do not have so much understanding that they can be entrusted with so much authority in their own affairs. But one ought not, when the Free Church has this principle, to assert that it does not have it.

It is self-evident that when the Free Church meetings cannot command and prescribe anything, then they cannot forbid the congregations anything either.

But then there comes the question, which is put to us with full justification: Why in the world does one have Free Church meetings, if these meetings have no authority either to command or to forbid the congregations anything?

Yes, it is an altogether legitimate question, one to which all Free Church people must be prepared to give an answer.

There is no one who is so fond of voluntary meetings as the Free Church people, and they must surely have a reason for it. There must be something that draws them together; otherwise they would not come in such numbers.

The purpose of these meetings, which work so strengthening and enlivening an effect both upon those who are guests and upon those who receive the guests, is, in its general character, to encourage one another to love and good works.

That those who carry on association politics in the old style cannot understand what such a thing is for is only too natural. Those, on the other hand, who are beginning to glimpse something, however little, of the way of the Spirit and the nature of freedom set great value on the most brotherly urging and counsel, even though no compulsion can be laid upon the congregations. They know what Paul calls exhortation in Christ, love’s encouragement, and the fellowship of the Spirit, tender love and mercies. And experience demonstrates that free Christian meetings are serviceable precisely for this.

Now it is entirely true that these are lofty and difficult things, and that it is very easy for the Free Church meetings to strike into the political track of the association meetings....

What is at issue is the unity of the Spirit, not the compulsion of the Law. If Spirit and love are not present, then let us go back to the old association order; for then it is only self-deception if we imagine that we have come out of it.

Georg Sverdrup
